\section*{Заключение}
\addcontentsline{toc}{section}{Заключение}

В результате проведенного исследования подтверждена гипотеза о том, что 
автоматизация логистических процессов на базе специализированного веб-приложения 
позволяет существенно повысить эффективность работы служб экспресс-доставки. 
Теоретический анализ показал, что классическая задача маршрутизации транспортных 
средств (VRP) в современных условиях требует адаптации под динамические изменения 
городского трафика и специфику «последней мили».

На основе полученных при разработке данных можно сформулировать следующие 
\textbf{выводы}:

\begin{enumerate}
    \item Использование открытых ГИС-технологий и компонентного подхода (React/Node.js) 
    обеспечивает необходимую архитектурную гибкость, позволяя интегрировать алгоритмы 
    оптимизации без зависимости от проприетарных платных сервисов.
    \item Проектирование модели данных с поддержкой пространственных расширений 
    (PostGIS) является критическим условием для масштабируемости системы, так как 
    позволяет эффективно обрабатывать геозапросы в режиме реального времени.
    \item Внедрение эвристических методов построения маршрутов демонстрирует 
    устойчивый экономический эффект, выраженный в снижении совокупного пробега 
    автотранспорта и сокращении времени планирования рейсов диспетчером.
\end{enumerate}

\textbf{Рекомендации и предложения.} Для практического применения разработанного 
решения в деятельности малых логистических компаний предлагается использовать 
модель облачного развертывания. Это обеспечит бесперебойный доступ к интерфейсу 
управления и мобильному клиенту курьера при минимальных затратах на ИТ-инфраструктуру.

\textbf{Перспективы углубления темы.} В дальнейшем исследовании планируется 
развитие функциональности системы за счет внедрения методов предиктивной аналитики 
для прогнозирования времени прибытия (ETA) с учетом исторических данных о заторах, 
а также разработки модуля динамической корректировки маршрута при поступлении 
срочных заявок в процессе выполнения рейса.


При разработке пояснительной записки к курсовому проекту строго соблюдались 
требования~\cite{gost_report}, определяющие структуру и правила оформления 
научно-исследовательских работ.

